\scanpage{21}
\mspage{8}

et b) $\forall i_0\in I$, \inserted{\unsure{nos}} $\alpha_{i_0},\beta_{i_0}\in C$ dans le diagramme \unsure{de|ci} p.~1.

\deleted{\uline{\unsure{Corollaire}}} De plus (supposant donné $W$ comme dans prop.~1 cor.~1), \inserted{les conditions suivantes} \deleted{on a $u_\bullet\in W_{I^\circ}$ (i.e. $u_{i_0}\in W$ $\forall i_0\in I$) alors $\unsure{u_I}\in TC=C\cap W$ et pour \deleted{\unsure{tout}} $i_0\in I$, on a $\widetilde u_{i_0}\in TC$ \unsure{donc}} sur $u_\bullet$ entre $X_\bullet,Y_\bullet$ $C$-cofibrants sont équivalentes

\[
\begin{aligned}
\deleted{\unsure{2}}1^\circ)&\quad u_\bullet\in(TC)_I,\quad\text{i.e. les }\widetilde u_{i_0}\in TC\quad(\forall i_0\in I)\\
\deleted{\unsure{3}}2^\circ)&\quad u_\bullet\in T_I\cap W_{I^\circ},\quad\text{i.e. }\forall i_0\in I,\quad\widetilde u_{i_0}\in C,\ u_{i_0}\in W,
\end{aligned}
\]

\deleted{\unsure{Alors}} $3^\circ)$ Pour tout $i_0\in I$, on a

\[
 \uline{\alpha_{i_0}},\beta_{i_0},\uline{\widetilde u_{i_0}},u_{i_0}\in TC
\]

(i.e. $\alpha_{i_0}$ \unsure{\&} $\widetilde u_{i_0}\in TC$, ce qui implique $\beta_{i_0}\in TC$ puis \inserted{$u_{i_0}\in TC$} par \unsure{cochgt.|cochgt$^{\rm t}$.} de base et par comp.) et ces conditions impliquent \inserted{$4^\circ)$} que $u_I:X_I\to Y_I$ est dans $TC$.

Voyons comment \deleted{\unsure{ou dans|on dé}} quelle mesure \unsure{le|ce} corollaire se déduit de la proposition. On a bien sûr

\diagram{\begin{tikzcd}[ampersand replacement=\&,column sep=small,row sep=small]
 \&3^\circ\arrow[dl,Rightarrow]\arrow[dr,Rightarrow]\&\\
 1^\circ\&\&2^\circ
\end{tikzcd}}

tautologiquement. D’ailleurs, la prop. appliquée à $TC$ \unsure{au lieu de} $C$ montre qu’on a $1^\circ\Rightarrow 3^\circ\mathbin{\&}4^\circ$, car en l’appliquant aux ens $I_{<i_0}$ on trouve que $\alpha_{i_0}\in TC$ $\forall i_0\in I$. Donc la proposition 2 implique

\[
 (1^\circ\Longleftrightarrow3^\circ)\Longrightarrow 2^\circ(\mathbin{\&}4^\circ).
\]
% archive sheet number: 20

